\chapter{Exécution de bout en bout}
\label{chap:execution-bout-en-bout}

\section{Commande cliente et décision de reconstitution}

La chaîne commence par une demande externe. Dans l'exécution étudiée, \Agent{CMD\_1} porte une commande de 20 unités alors que le stock initial de produit fini est nul. Le modèle conserve cette commande en attente et laisse la politique autonome évaluer le besoin de reconstitution.

Cette évaluation crée \Agent{REAPPRO\_1}, un ordre interne de 20 unités. Les deux objets n'ont pas la même fonction. \Agent{CMD\_1} représente l'obligation de service envers le client; \Agent{REAPPRO\_1} autorise le rétablissement du stock nécessaire à ce service. L'ordre interne n'est ni une seconde commande externe ni une livraison à compter dans les KPI clients.

Le passage par la politique de stock préserve une logique Make-to-Stock. Make n'est pas déclenché par la conversion directe de la demande en ordre de fabrication. Il répond à une décision interne fondée sur l'état du stock. Une fois le lot achevé, le stock physique est crédité et la commande en attente peut être réévaluée.

\section{Analyse matière et intervention de Source}

La nomenclature de \Agent{REAPPRO\_1} exige 250 unités de GPL, 20 bouteilles et 20 kits d'accessoires. Les bouteilles et les accessoires sont disponibles au début du scénario, contrairement au GPL. Le diagnostic matière impose donc Source avant Make.

La réception de 250 unités de GPL crédite le stock matière, puis l'événement \texttt{MaterialAvailable} informe Make que la contrainte amont est levée. La quantité de GPL traverse plusieurs opérations Source; elle ne doit pas être additionnée à chaque ligne de la trace. Le flux physique et le message de disponibilité décrivent deux aspects complémentaires du même passage, sans être interchangeables.

Les délais Source et Deliver utilisent une conversion temporelle, tandis que Make conserve sa dynamique propre. Les timestamps de la trace ordonnent les événements; ils ne suffisent pas à reconstruire les durées métier exportées par le pipeline VSM.

\section{Make et crédit du stock fini}

Make débute après la disponibilité matière et produit les 20 unités de l'ordre interne. Le transfert M1.5 vers le magasin constitue l'étape où le stock physique de produit fini est crédité. Les jetons réservés à l'animation sont exclus de ce comptage.

La fin d'une opération de production ne rend pas automatiquement la commande livrable. Le modèle clôt d'abord l'ordre interne, puis relance l'analyse des commandes en attente. \Agent{CMD\_1} retrouve alors un stock disponible dont la provenance reste liée à \Agent{REAPPRO\_1}. Cette relation est exportée dans les deux sens dans l'ABox.

Les mécanismes de rebut et de reprise sont disponibles dans le modèle. Le scénario retenu ne comporte toutefois pas de défauts permettant d'évaluer quantitativement leur comportement.

\section{Réveil de la commande et Deliver}

Après le crédit du stock, la politique réveille \Agent{CMD\_1}. Deliver réserve les 20 unités, prépare l'expédition et assure le transport jusqu'au client. Le stock est décrémenté lors de la satisfaction effective de la commande. L'exécution se termine avec 20 unités commandées, 20 unités livrées et un stock fini nul.

La commande est complète, mais tardive. Ce résultat porte sur une seule demande et ne représente pas un taux de service de long terme. Il confirme en revanche que le service client et la reconstitution interne restent séparés dans les comptes.

\FigureOrPlaceholder{figures/chaine_execution_consolidee.png}{Chaîne fonctionnelle de satisfaction d'une commande lorsque les stocks imposent une reconstitution préalable.}{fig:chaine-execution-consolidee}

\section{Portée de la trace}

La chronologie vérifiée est regroupée dans la \cref{tab:chronologie-execution-vsm}. Elle établit l'ordre entre diagnostic, Source, Make, crédit du stock, réveil de la commande, Deliver et clôture. Les valeurs VSM utilisent un autre niveau d'instrumentation et sont analysées au chapitre~\ref{chap:analyse-vsm}.

La trace et l'ABox établissent ensemble la présence de \Agent{CMD\_1} et \Agent{REAPPRO\_1}, leur dépendance, l'approvisionnement GPL, la livraison complète et le stock final nul. Elles ne quantifient pas l'effet propre de la politique de stock ou du retard fournisseur. Une comparaison contrôlée serait nécessaire pour isoler ces contributions.
